\documentclass[11pt]{article}
\usepackage[utf8]{inputenc}
\usepackage[spanish]{babel}
\usepackage{amsmath,amsthm,amsfonts,amssymb,amscd}
\usepackage{multirow,booktabs}
\usepackage[table]{xcolor}
\usepackage{fullpage}
\usepackage{lastpage}
\usepackage{enumitem}
\usepackage{fancyhdr}
\usepackage{mathrsfs}
\usepackage{wrapfig}
\usepackage{setspace}
\usepackage{calc}
\usepackage{multicol}
\usepackage[retainorgcmds]{IEEEtrantools}
\usepackage[margin=3cm]{geometry}
\usepackage{empheq}
\usepackage{framed}
\usepackage[most]{tcolorbox}
\usepackage{soul}
\usepackage{longtable}
\usepackage{array}
\usepackage{ragged2e}
\newlength{\tabcont}
\setlength{\parindent}{0.0in}
\setlength{\parskip}{0.05in}
\colorlet{shadecolor}{orange!15}
\parindent 0in
\parskip 12pt
\geometry{margin=1in, headsep=0.25in}

\newcolumntype{L}[1]{>{\RaggedRight\arraybackslash}p{#1}}

\begin{document}
\setcounter{section}{0}

\thispagestyle{empty}

\begin{center}
{\LARGE \bf Análisis de Riesgos}\\[4pt]
{\large Tesis I}\\[2pt]
Roberto Treviño Cervantes
\end{center}

\section{Contexto}

La presente tesis aborda el diseño de perfiles de trayectoria de cuarto orden (\textit{snap-limited}) optimizados por enjambre de partículas (PSO), utilizando una red neuronal convolucional (CNN) entrenada con el conjunto WM-811K como función de costo dinámica. La validación se realiza íntegramente sobre una simulación del multi-cuerpo de Al-Rawashdeh (2022). Bajo este alcance —exclusivamente simulación numérica— los riesgos relevantes son aquellos asociados a la disponibilidad de datos, la convergencia del entrenamiento y la capacidad de cómputo, no al acceso a equipo físico de litografía.

\section{Matriz de Riesgos}

La Tabla~\ref{tab:riesgos} sintetiza los riesgos identificados, su impacto sobre los entregables del proyecto y la estrategia de mitigación correspondiente.

\begingroup
\small
\setlength{\tabcolsep}{6pt}
\renewcommand{\arraystretch}{1.5}
\begin{longtable}{L{3.8cm} L{2.2cm} L{6.5cm}}
\caption{Matriz de riesgos del proyecto de tesis.}
\label{tab:riesgos} \\
\toprule
\textbf{Riesgo} & \textbf{Impacto} & \textbf{Mitigación} \\
\midrule
\endhead
\bottomrule
\endlastfoot

No obtener datos &
Alto &
El entrenamiento de la CNN depende del acceso al conjunto WM-811K. Si las condiciones de licencia se vuelven restrictivas, se genera un conjunto sintético equivalente a partir de perturbaciones controladas de trayectoria sobre el gemelo digital, replicando las clases de defecto morfológico de WM-811K. \\[4pt]

Modelo no converge &
Medio &
La función de costo aprendida no es convexa y el acoplamiento CNN--PSO puede estancarse en mínimos locales. Se simplifica el problema reduciendo la dimensión del espacio de búsqueda $(v_{max},\, a_{max},\, j_{max},\, s_{max})$, fijando subconjuntos de parámetros y empleando una inicialización estrategia antes de escalar al problema completo. \\[4pt]

Hardware no disponible &
Alto &
El riesgo se refiere a la capacidad de cómputo requerida para entrenar la CNN y ejecutar PSO sobre la simulación. Se puede utilizar la configuración reducida sin rotaciones como banco de pruebas de menor escala, conservando la validez metodológica del estudio. \\[4pt]

\end{longtable}
\endgroup

\end{document}
